% Options for packages loaded elsewhere
\PassOptionsToPackage{unicode}{hyperref}
\PassOptionsToPackage{hyphens}{url}
\documentclass[
  a4paper,
]{ltjsarticle}
\usepackage{xcolor}
\usepackage[margin=25mm]{geometry}
\usepackage{amsmath,amssymb}
\setcounter{secnumdepth}{-\maxdimen} % remove section numbering
\usepackage{iftex}
\ifPDFTeX
  \usepackage[T1]{fontenc}
  \usepackage[utf8]{inputenc}
  \usepackage{textcomp} % provide euro and other symbols
\else % if luatex or xetex
  \usepackage{unicode-math} % this also loads fontspec
  \defaultfontfeatures{Scale=MatchLowercase}
  \defaultfontfeatures[\rmfamily]{Ligatures=TeX,Scale=1}
\fi
\usepackage{lmodern}
\ifPDFTeX\else
  % xetex/luatex font selection
\fi
% Use upquote if available, for straight quotes in verbatim environments
\IfFileExists{upquote.sty}{\usepackage{upquote}}{}
\IfFileExists{microtype.sty}{% use microtype if available
  \usepackage[]{microtype}
  \UseMicrotypeSet[protrusion]{basicmath} % disable protrusion for tt fonts
}{}
\makeatletter
\@ifundefined{KOMAClassName}{% if non-KOMA class
  \IfFileExists{parskip.sty}{%
    \usepackage{parskip}
  }{% else
    \setlength{\parindent}{0pt}
    \setlength{\parskip}{6pt plus 2pt minus 1pt}}
}{% if KOMA class
  \KOMAoptions{parskip=half}}
\makeatother
\usepackage{longtable,booktabs,array}
\newcounter{none} % for unnumbered tables
\usepackage{calc} % for calculating minipage widths
% Correct order of tables after \paragraph or \subparagraph
\usepackage{etoolbox}
\makeatletter
\patchcmd\longtable{\par}{\if@noskipsec\mbox{}\fi\par}{}{}
\makeatother
% Allow footnotes in longtable head/foot
\IfFileExists{footnotehyper.sty}{\usepackage{footnotehyper}}{\usepackage{footnote}}
\makesavenoteenv{longtable}
\setlength{\emergencystretch}{3em} % prevent overfull lines
\providecommand{\tightlist}{%
  \setlength{\itemsep}{0pt}\setlength{\parskip}{0pt}}
% Glyph map for lualatex (newunicodechar). Maps symbols that may lack font glyphs.
\usepackage{newunicodechar}
\usepackage{amssymb}
\newunicodechar{∎}{\ensuremath{\blacksquare}}
\newunicodechar{✓}{\ensuremath{\checkmark}}
\newunicodechar{✗}{\ensuremath{\times}}
\newunicodechar{⟹}{\ensuremath{\Longrightarrow}}
\newunicodechar{⟺}{\ensuremath{\Longleftrightarrow}}
\newunicodechar{↪}{\ensuremath{\hookrightarrow}}
\newunicodechar{⌊}{\ensuremath{\lfloor}}
\newunicodechar{⌋}{\ensuremath{\rfloor}}
\newunicodechar{≃}{\ensuremath{\simeq}}
\newunicodechar{≈}{\ensuremath{\approx}}
\newunicodechar{≤}{\ensuremath{\leq}}
\newunicodechar{≥}{\ensuremath{\geq}}
\newunicodechar{≠}{\ensuremath{\neq}}
\newunicodechar{→}{\ensuremath{\rightarrow}}
\newunicodechar{←}{\ensuremath{\leftarrow}}
\newunicodechar{↔}{\ensuremath{\leftrightarrow}}
\newunicodechar{⊥}{\ensuremath{\perp}}
\newunicodechar{√}{\ensuremath{\surd}}
\newunicodechar{∑}{\ensuremath{\sum}}
\newunicodechar{∏}{\ensuremath{\prod}}
\newunicodechar{∞}{\ensuremath{\infty}}
\newunicodechar{∈}{\ensuremath{\in}}
\newunicodechar{∀}{\ensuremath{\forall}}
\newunicodechar{∣}{\ensuremath{\mid}}
\newunicodechar{γ}{\ensuremath{\gamma}}
\newunicodechar{λ}{\ensuremath{\lambda}}
\newunicodechar{μ}{\ensuremath{\mu}}
\newunicodechar{ρ}{\ensuremath{\rho}}
\newunicodechar{σ}{\ensuremath{\sigma}}
\newunicodechar{φ}{\ensuremath{\varphi}}
\newunicodechar{ψ}{\ensuremath{\psi}}
\newunicodechar{θ}{\ensuremath{\theta}}
\newunicodechar{Δ}{\ensuremath{\Delta}}
\newunicodechar{Π}{\ensuremath{\Pi}}
\newunicodechar{Λ}{\ensuremath{\Lambda}}
\newunicodechar{τ}{\ensuremath{\tau}}
\newunicodechar{ε}{\ensuremath{\varepsilon}}
\newunicodechar{ω}{\ensuremath{\omega}}
\newunicodechar{π}{\ensuremath{\pi}}
\newunicodechar{ν}{\ensuremath{\nu}}
\newunicodechar{±}{\ensuremath{\pm}}
\newunicodechar{×}{\ensuremath{\times}}
\newunicodechar{·}{\ensuremath{\cdot}}
\usepackage{bookmark}
\IfFileExists{xurl.sty}{\usepackage{xurl}}{} % add URL line breaks if available
\urlstyle{same}
\hypersetup{
  pdftitle={高対称床のゼロ閉塞因数分解------偶数系 N=8k の \textbackslash pm i 光子対完全被覆と奇数系の最小 p 波凝縮体},
  pdfauthor={Noriaki Kihara (WF System Co., Ltd.), ORCID 0009-0004-6753-4020},
  hidelinks,
  pdfcreator={LaTeX via pandoc}}

\title{高対称床のゼロ閉塞因数分解------偶数系 \(N=8k\) の \(\pm i\)
光子対完全被覆と奇数系の最小 \(p\) 波凝縮体}
\author{Noriaki Kihara (WF System Co., Ltd.), ORCID 0009-0004-6753-4020}
\date{2026-09-12}

\begin{document}
\maketitle

\textbf{シリーズ:} 自己無撞着な関係波閉鎖系の基礎（論文4／全10本）\\
\textbf{著者:} 木原 範昭（WF System Co., Ltd.）　\textbf{日付:}
2026-09-12\\
\textbf{Version DOI:} 10.5281/zenodo.22729081\\
\textbf{Concept DOI:} 10.5281/zenodo.22729080\\
\textbf{ORCID:} 0009-0004-6753-4020\\
\textbf{Zenodo:} https://zenodo.org/records/22729081\\

\begin{center}\rule{0.5\linewidth}{0.5pt}\end{center}

\subsection{要旨}\label{ux8981ux65e8}

高対称床の二乗閉塞 \(\sum_e z_e^2=0\)
を、より小さなゼロ閉塞ブロックへの因数分解として調べ、最小ブロックサイズが
\(N\) の数論で決まる\textbf{巡回二乗波の因数分解定理}に一本化する。巡回
Fourier 相対平衡床 \(z_{ij}=r_d e^{i2\pi(i+j)/N}\)
を定義式から再構成し、全 unordered
pair・3波・4波を総当たり監査した。\textbf{定理}：固定距離クラス内の二乗波
\(w=z^2\) は位相ステップ \(q=e^{i8\pi/N}\)（位数 \(L=N/\gcd(N,4)\)）で正
\(L\) 角形を巡回し、最小ゼロ閉塞ブロックは \(L\) の最小素因子
\(p=\operatorname{spf}(L)\) 角形。ゆえ \(p=2\iff L\) 偶
\(\iff 8\mid N\)------\textbf{同一距離クラス内で厳密な二波自己閉塞
\(z_b=\pm i\,z_a\)（光子型 \(\pm i\) 対）が可能なのは \(N=8k\)
に限る（任意 \(N\)
で証明）}。全空間（異距離クラス間を含む）でも同じ選択則を全
\(N=3\)--\(40\) の全 pair 監査で確認した（位相格子＋パリティ＋振幅反射
\(r_d=r_{N/2-d}\) による；反射以外の偶然の等振幅対がない任意 \(N\)
一般証明は今後の課題、§8）。\(N=8k\) では \(m=N/8\) の
\(z_{i+m}=\pm i\,z_i\) によりグラフ探索なしに全 \(M\) 波を \(M/2\)
個の光子対へ\textbf{明示的に完全 matching}できる（候補辺数の閉形式
\(E_\gamma=N(4N-11)/2=84,424,1020,1872,2980\)、候補グラフは二部・次数
\(2/4/8\)・連結成分サイズ \(4/8/16\)、matching は非一意）。\(N=8k\)
以外の偶数 \(N\) は最小 \(p\) 角形ゼロ閉塞へ100\%分解でき、各距離クラス
\(C_d\) 自身が独立に \(\sum_{e\in C_d}z_e^2=0\)。奇数 \(N\)
は2波対0で、最小素因子 \(p\) 波（\(3\mid N\) は3波、奇素数 \(N\) は
\(N\)
波）が基本単位。因数分解は一般に非一意。床上ではこれら全ブロックが単一
carrier に位相ロックされ周波数比は \(1{:}1\)（代数的に分解可能 ≠
力学的に独立）。分類：静的な代数的因数分解の存在（粒子同定・力学保存は別課題）。

\begin{center}\rule{0.5\linewidth}{0.5pt}\end{center}

\subsection{1.
課題------閉塞の因数分解}\label{ux8ab2ux984cux9589ux585eux306eux56e0ux6570ux5206ux89e3}

論文1・3の床は
\(\sum_e z_e^2=0\)（零閉塞）を満たす。この総閉塞が、より小さなゼロ閉塞ブロックの直和に\textbf{因数分解}できるか、そのブロックの型が
\(N\)
でどう変わるかを調べる。最小の非自明ブロックは2波（\(z_2=\pm i z_1\)、光子型
\(\pm i\) 対）である（論文9）。

\subsection{2. 対象と方法}\label{ux5bfeux8c61ux3068ux65b9ux6cd5}

高対称理論床 \(z_{ij}=r_d e^{i2\pi(i+j)/N}\)（論文2の(C)、巡回 Fourier
相対平衡）を定義式から再構成。全空間で
\(\|z\|^2\)、\(\sum z\)、\(\sum z^2\)、固有残差を独立検算。二波条件を無次元残差
\(\varepsilon_{ij}=|z_i^2+z_j^2|/(|z_i|^2+|z_j|^2)\) で全 pair
に総当たり（strict は \(\varepsilon\le10^{-12}\)）、3波・4波は
meet-in-the-middle
で総当たり。\(Z_0=(v+\delta g)/\|\cdot\|\)、\(\delta=10^{-15}\) でも
strict 判定の \(N\)・個数は不変。

\subsection{\texorpdfstring{3.
巡回二乗波の因数分解定理（\(L=N/\gcd(N,4)\)）}{3. 巡回二乗波の因数分解定理（L=N/\textbackslash gcd(N,4)）}}\label{ux5de1ux56deux4e8cux4e57ux6ce2ux306eux56e0ux6570ux5206ux89e3ux5b9aux7406lngcdn4}

固定距離クラス \(C_d\) の辺 \(\{i,i+d\}\) は
\(z_{i,i+d}=r_d e^{i2\pi(2i+d)/N}\) ゆえ、二乗波

\[w_i=z_{i,i+d}^2=r_d^2\,e^{i4\pi(2i+d)/N},\qquad w_{i+1}=q\,w_i,\quad q=e^{i8\pi/N},\quad L=\operatorname{ord}(q)=\frac{N}{\gcd(N,4)}.\]

クラス内の二乗波は正 \(L\) 角形の頂点を巡回するので、\(q^{L/p}\) が位数
\(p=\operatorname{spf}(L)\) の部分群＝正 \(p\) 角形を張り、頂点和
\(=0\)------\textbf{\(p\) 波零閉塞ブロックが存在する}。さらに
\textbf{Lam--Leung の消失和定理［2］}より、\(L\)
乗根の非空な消失和の最小重み（項数）は \(L\)
の最小素因子に等しい。距離クラス内の二乗波は等モジュラス \(r_d^2\) の
\(L\) 乗根ゆえ、この \(p\)
角形は単なる構成例でなく\textbf{同一クラス内の最小零閉塞ブロック}である（存在だけでなく最小性まで）。これで偶数・奇数を同じ数式が説明する：

{\def\LTcaptype{none} % do not increment counter
\begin{longtable}[]{@{}
  >{\raggedright\arraybackslash}p{(\linewidth - 4\tabcolsep) * \real{0.3333}}
  >{\raggedright\arraybackslash}p{(\linewidth - 4\tabcolsep) * \real{0.3333}}
  >{\raggedright\arraybackslash}p{(\linewidth - 4\tabcolsep) * \real{0.3333}}@{}}
\toprule\noalign{}
\begin{minipage}[b]{\linewidth}\raggedright
\(N\)
\end{minipage} & \begin{minipage}[b]{\linewidth}\raggedright
\(L=N/\gcd(N,4)\)
\end{minipage} & \begin{minipage}[b]{\linewidth}\raggedright
最小ブロック \(p=\operatorname{spf}(L)\)
\end{minipage} \\
\midrule\noalign{}
\endhead
\bottomrule\noalign{}
\endlastfoot
奇数 & \(N\) &
\(\operatorname{spf}(N)\)（\(3\mid N\)→3波、\(N=25,35\)→5波、奇素数
\(N\)→\(N\)波） \\
\(N\equiv2\ (\mathrm{mod}\,4)\) & \(N/2\) &
\(\operatorname{spf}(N/2)\) \\
\(N\equiv4\ (\mathrm{mod}\,8)\) & \(N/4\)（奇） &
\(\operatorname{spf}(N/4)\) \\
\(N\equiv0\ (\mathrm{mod}\,8)\) & \(N/4\)（偶） & \(2\)（\(\pm i\)
光子対） \\
\end{longtable}
}

ゆえに
\(\boxed{p=2\iff L\ \text{偶}\iff 8\mid N}\)------\textbf{同一距離クラス内で2波零閉塞（\(\pm i\)
光子対）が可能なのは \(8\mid N\) に限る}。\(\operatorname{ord}(q)=L\) と
\(p\) 角形和 \(=0\) は全 \(N=3\)--\(40\) で確認（\(N=4\) は \(L=1\)
の退化で本機構の対象外＝別機構で3波閉塞、§6）。

\textbf{なぜ「8」か（2進的選択則）。} 元の巡回位相は \(i\to i+1\) で
\(\Delta\varphi=4\pi/N\)、二乗写像でさらに2倍され
\(\Delta\arg(z^2)=8\pi/N\)。二波零閉塞 \(w_b=-w_a\) には二乗波の巡回群
\(\mathbb Z_L\) に対蹠点（位数2の部分群）が要り、それは \(L\) 偶
\(\iff\nu_2(N)\ge3\iff 8\mid N\)（\(\nu_2\) は2進付値）。すなわち
\(N=8k\)
は偶然の周期ではなく、\textbf{二乗写像が要求する2進的選択則}である。

\textbf{\(N=8k\) の明示構成（グラフ探索不要）。} \(L=N/4=2k\)、\(m=N/8\)
とすると \(q^m=e^{i\pi}=-1\)。ゆえ同一クラスで

\[w_{i+m}=-w_i\ \Longrightarrow\ z_{i+m}^2=-z_i^2\ \Longrightarrow\ z_{i+m}=\pm i\,z_i\]

（振幅 \(r_d\) 一致・位相差ちょうど \(\pi/2\)）。各距離クラスで写像
\(i\mapsto i+m\) の対蹠点（\(w_{i+m}=-w_i\)）を組にすれば完全 matching
がその場で作れる（\(i\mapsto i+m\) の軌道長は通常クラス \(d<N/2\)
が8、直径クラス \(d=N/2\)
が4、いずれも偶ゆえ軌道を過不足なく対に分割できる）。\(N=8\)--\(40\) の
\(8k\) で厳密 \(\pm i\)・全 \(M\) 辺を \(M/2\)
対に被覆を確認（\(14,60,138,248,390\)）。§4 の perfect matching
はこの\textbf{明示的構成}を持つ。

\textbf{異距離クラス間を含めても \(\iff 8\mid N\)。} \(\pm i\) 対は
\(|z_a|=|z_b|\) かつ位相差ちょうど \(\pi/2\) を要する。(1) 床の位相は
\((2\pi/N)\mathbb Z\) 上ゆえ \(\pi/2\) 位相差は \(4\mid N\)
が必要（\(4\nmid N\)＝奇数・\(N\equiv2\ (\mathrm{mod}\,4)\)
では不可）。(2) 位相 \(2\pi(2i+d)/N\) の \(2i\) は偶ゆえ、クラス
\(d,d'\) 間で \(\pi/2\) には \(d'-d\equiv N/4\ (\mathrm{mod}\,2)\)
が必要。(3) 振幅が等しいのは反射 \(r_d=r_{N/2-d}\)（\(N=3\)--\(40\)
で数値確認、誤差 \(\le1.2\times10^{-16}\)）で
\(d'-d\equiv N/2\ (\mathrm{mod}\,2)\)：\(N\equiv4\ (\mathrm{mod}\,8)\)
は \(N/2\) 偶だが \(N/4\) 奇で両立せず、\(8\mid N\)
のみ実現する。全ペア（異クラス込み）の厳密 \(\pm i\)
対数が既往の網羅監査（§4）と独立に完全一致することを確認した。以下、この定理の帰結として偶数系（§4）・奇数系（§5）・距離クラス独立閉塞（§6）を述べる。

\subsection{\texorpdfstring{4. 偶数系------\(N=8k\) の \(\pm i\)
光子対完全被覆}{4. 偶数系------N=8k の \textbackslash pm i 光子対完全被覆}}\label{ux5076ux6570ux7cfbn8k-ux306e-pm-i-ux5149ux5b50ux5bfeux5b8cux5168ux88abux8986}

\begin{itemize}
\tightlist
\item
  \textbf{厳密 \(z_b=\pm i z_a\) が存在するのは
  \(N=8k\)（8,16,24,32,40）のみ}（他の偶数 \(N\) は候補0）------§3
  定理の \(p=2\iff8\mid N\) の帰結。以下はその組合せ構造。
\item
  全 \(M\) 波を重複なしの \(M/2\) 個の二波自己閉塞対へ\textbf{完全
  matching}可能：\(N=8{:}14/28\)、\(16{:}60/120\)、\(24{:}138/276\)、\(32{:}248/496\)、\(40{:}390/780\)。
\item
  候補辺数の閉形式
  \(E_\gamma=N(4N-11)/2=84,424,1020,1872,2980\)。候補グラフは全 \(N=8k\)
  で\textbf{二部}、次数は \(2/4/8\) のみ（\(\#\deg2=N/2\),
  \(\#\deg4=N\), \(\#\deg8=M-3N/2\)）、連結成分サイズは \(4/8/16\)
  のみ（\(C_4=N/8\), \(C_8=N/8\), \(C_{16}=N(N-4)/32\)）。
\item
  \textbf{選択則（parity）}：位相 \(90^\circ\) には \(\Delta(i+j)=N/4\)
  が必要で、距離 \(d\) の辺は
  \(i+j\equiv d\ (\mathrm{mod}\,2)\)。\(N=8k\) なら \(N/4\) 偶数で
  parity 保存、\(N=8k+4\) なら奇数で反転し等振幅クラスと両立不能。
\item
  \textbf{完全 pair 分解は非一意}：perfect matching
  は元の辺を共有する対を0組にも最大 \(N\)
  組にも選べる。静的状態だけからは「どの2波が一個の光子か」は一意でなく、\(+i/-i\)（右旋/左旋）の絶対個数も無向
  pair だけからは定義できない（伝播方向・orientation が別途必要）。
\end{itemize}

\subsection{\texorpdfstring{5. 奇数系------最小 \(p\)
波凝縮体}{5. 奇数系------最小 p 波凝縮体}}\label{ux5947ux6570ux7cfbux6700ux5c0f-p-ux6ce2ux51ddux7e2eux4f53}

\begin{itemize}
\tightlist
\item
  奇数 \(N\)（3,5,\ldots,39）では\textbf{光子型二波対が完全に0}（§3
  定理より奇数は \(L=N\)・\(p=\operatorname{spf}(N)\ge3\)）、最大
  matching 0、全 \(M\) 波が pair 外。それでも最小 \(p\)
  波ゼロ閉塞凝縮体へ\textbf{100\%完全分解}可能（constructive coverage
  \(=1.0\)、最大群残差 \(\le1.2\times10^{-15}\)）。
\item
  \(z^2\) が \(N\) 乗根を一巡し、\(p=\text{smallest prime factor}(N)\)
  の正 \(p\) 角形へ分割：\(3\mid N\to3\)波、\(N=5,25,35\to\)
  最小5波、奇素数 \(N\to\) 最小 \(N\) 波。
\item
  全3波総当たりで exact triad が出るのは
  \(3\mid N\)（\(N=3,9,15,21,27,33,39\)）だけ。全4波は全奇数系で0。
\end{itemize}

\subsection{6.
残余波・距離クラスの独立閉塞（偶数系）}\label{ux6b8bux4f59ux6ce2ux8dddux96e2ux30afux30e9ux30b9ux306eux72ecux7acbux9589ux585eux5076ux6570ux7cfb}

\begin{itemize}
\tightlist
\item
  \(N=8k\) は残余0（完全 pair 被覆）、それ以外の偶数 \(N\) は exact
  2波対0で残余＝全 \(M\)
  波。\textbf{残余全体も機械精度でゼロ閉塞}（相対残差
  \(4.1\times10^{-17}\)--\(2.4\times10^{-16}\)）。
\item
  \(N>4\) では\textbf{各距離クラス \(C_d\) 自身が独立に
  \(\sum_{e\in C_d}z_e^2=0\)}。総閉塞は距離クラス凝縮体へ既に分解される。二乗波
  \(w=z^2\) の同一クラス内位相ステップ \(q=e^{i8\pi/N}\)、位数
  \(L=N/\gcd(N,4)\)。\(L\) の因数 \(p\) ごとに \(p\)
  角形ゼロ閉塞へ分解可能。
\item
  全偶数 \(N\) で全 \(M\)
  波を重複なしの最小対称ゼロ閉塞ブロックへ100\%分解（例：\(N=4{:}3\times2\),
  \(6{:}3\times5\), \(8{:}2\times14\), \(10{:}5\times9\),
  \(12{:}3\times22\), \ldots, \(40{:}2\times390\)）。exact 3波閉塞は
  \(N=4,6,12,18,24,30,36\)。因数分解は \(L\)
  が合成数なら非一意（\(N=24{:}2/3/6\), \(30{:}3/5/15\),
  \(40{:}2/5/10\)）。
\end{itemize}

\subsection{7. 床上では単一
carrier------倍音なし}\label{ux5e8aux4e0aux3067ux306fux5358ux4e00-carrierux500dux97f3ux306aux3057}

本論文が因数分解した巡回Fourier床（§2、\(z_{ij}=r_d e^{i2\pi(i+j)/N}\)）そのものの凍結生成子
\(K\) の \(iK\) 正固有値を監査（監査床は §2--6
と同一・\texttt{build\_theoretical\_floor} で
\(z=r_d e^{i\theta}\)、\(\theta=2\pi(i+j)/N\)
から再構成。90°整数K床＝論文3 とは別メンバー）：

\begin{itemize}
\tightlist
\item
  高対称床の占有状態は全波が一つの carrier
  を共巻きする\textbf{単一固有モード}（固有残差
  \(\le3.0\times10^{-14}\)）。因数分解しても床上のブロック間周波数比は
  \(1{:}1\)。
\item
  \(N=4\)--\(40\) 全偶数系で
  \textbf{2/3/4倍の整数倍音関係0件・分母16以下の低分母有理比0件・三波和共鳴
  \(\omega_i+\omega_j=\omega_k\) 0件}。\(N=8k\)
  の光子対完全被覆と倍音出現は無相関。
\item
  二波 \(z_j=\pm iz_i\) は \(90^\circ\) 位相差のゼロ閉塞幾何であって
  \(1{:}2\)
  等の倍音ではない。高対称背景は多数の閉塞凝縮体へ因数分解できるが、床上では独立時計を持たず同一
  carrier に位相ロックされている。
\end{itemize}

\textbf{対照（何が因数分解を作るか）。} 同じく自己無撞着・零閉塞
\(\sum z^2=0\) を満たし90°位相格子を持つ make\_parent
床について全二波を監査したが（実験15）、厳密な \(z_b=\pm i\,z_a\) 対は
\(N=3\)--\(40\) の全系列で \(\varepsilon\le10^{-9}\)
0件（\(\pm90^\circ\) 位相は自発形成されるが等振幅化が不足し、最良でも
\(\varepsilon=3.7\times10^{-8}\)）。したがって光子型二波因数分解は\textbf{零閉塞そのものや90°位相格子だけからは生じず}、本論文が扱う巡回Fourier床の位相巡回性（\(q=e^{i8\pi/N}\)）と距離クラス等振幅構造に依存する。この対照は、make\_parent
床と高対称理論床が同じ点火資格を満たす\textbf{別メンバー}であること（論文3・論文0）とも整合する。

\subsection{8.
主張分類・未決}\label{ux4e3bux5f35ux5206ux985eux672aux6c7a}

\begin{itemize}
\tightlist
\item
  分類：静的な代数的因数分解の\textbf{存在}。同一距離クラス内は位相のみで一般証明（§3
  定理 \(L=N/\gcd(N,4)\), \(p=\operatorname{spf}(L)\)）、\(N=8k\)
  明示構成は厳密、全空間の「\(\iff 8\mid N\)」は位相格子＋パリティ＋振幅反射で説明し全ペア監査（\(N=3\)--\(40\)）で数値確定。判定：保持。
\item
  未決：振幅の等値が反射 \(r_d=r_{N/2-d}\) に限る（\(d'-d\)
  が奇となる偶然の等振幅クラス対が存在しない）ことは \(N=3\)--\(40\)
  で数値確認。任意 \(N\) の一般証明（巡回Fourier相対平衡の \(r_d\)
  が反射以外に等値を持たないこと）は今後の検討課題とする。ただし
  \(4\nmid N\) の2波対不在は振幅に依らず位相格子から無条件に従う。
\item
  明示的に扱わないこと：粒子種の同定、力学的保存（固定 \(R_{124,23}\)
  交換散乱下でのブロック生成・消滅・再結合）は本論文の範囲外。\(\pm i\)
  対の helicity（右旋/左旋の絶対個数）は静的無向 pair
  だけからは定義できない。
\item
  対照：make\_parent 床では exact \(\pm i\) pair \(=0\)（全 pair
  監査・\(N=3\)--\(40\)、実験15）。因数分解は零閉塞・90°格子ではなく巡回Fourier床の追加構造に依存する（§7）。
\item
  高対称背景は「光子の海」より一般に\textbf{最小ゼロ閉塞凝縮体の海}で、\(N=8k\)
  は最小単位が2波（光子型）まで分解できる特別系列。情報の最小単位が2波である代数的根拠は論文9。
\end{itemize}

\subsection{関連研究（構造的一致・導出元でない）}\label{ux95a2ux9023ux7814ux7a76ux69cbux9020ux7684ux4e00ux81f4ux5c0eux51faux5143ux3067ux306aux3044}

\begin{itemize}
\tightlist
\item
  消失和（vanishing sums of roots of
  unity）の因数分解（Lam--Leung）：ゼロ閉塞ブロックが素サイズへ分解される構造と一致。
\item
  辺グラフ／association scheme（\(L(K_N)\), Johnson
  スキーム）：距離クラス構造。
\end{itemize}

\subsection{参考文献}\label{ux53c2ux8003ux6587ux732e}

\begin{enumerate}
\def\labelenumi{\arabic{enumi}.}
\tightlist
\item
  木原範昭「自己無撞着な関係波閉鎖系におけるインフレーション的急拡大の機構」Concept
  DOI 10.5281/zenodo.22112008.
\item
  T. Y. Lam and K. H. Leung, ``On vanishing sums of roots of unity'', J.
  Algebra 224 (2000) 91.
\item
  D. Cvetković, P. Rowlinson, S. Simić, \emph{An Introduction to the
  Theory of Graph Spectra}, Cambridge (2010).
\end{enumerate}

\subsection{再現}\label{ux518dux73fe}

監査スクリプト・CSVは各パッケージに一式：偶数系
\texttt{高対称理論床\_偶数系光子対精密解析\_20260908/}（\texttt{audit\_highsym\_even\_photon\_pairs\_v2.py}、\texttt{even\_N\_summary.csv}、\texttt{selected\_maximum\_matchings.csv}、\texttt{component\_size\_counts.csv}）、\texttt{高対称理論床\_偶数系\_光子対と凝縮体因数分解\_20260908/}（\texttt{analyze\_unpaired\_condensates\_v1.py}、\texttt{primitive\_zero\_closure\_groups.csv}、\texttt{distance\_class\_zero\_closure.csv}）、\texttt{高対称理論床\_偶数系\_倍音関係監査\_20260908/}（\texttt{analyze\_highsym\_harmonic\_spectrum\_v2.py}、\texttt{positive\_frequency\_modes.csv}）；奇数系
\texttt{高対称理論床\_奇数系\_ゼロ閉塞因数分解監査\_20260908/}（\texttt{analyze\_highsym\_odd\_closure\_factorization\_v1.py}、\texttt{constructive\_primitive\_zero\_closure\_groups.csv}）；make\_parent
床の光子対監査
\texttt{make\_parent型初期値\_自己無撞着構造\_20260907/二波自己閉塞光子対監査\_20260908/}。§3
の因数分解定理（\(L=N/\gcd(N,4)\)・\(p=\operatorname{spf}(L)\)・\(N=8k\)
明示構成・異クラス不在）の検証は
\texttt{巡回二乗波因数分解定理\_20260913/}（\texttt{verify\_cyclic\_square\_wave\_factorization\_20260913.py}＝位相定理＋\(N=8k\)
明示構成＋全ペア監査の再現の読出しのみ、\texttt{results/cyclic\_factorization\_summary.csv}・\texttt{実行ログ}、\texttt{分析}・README・SHA256SUMS・run\_all；入力は上記偶数系
CSV）。各に README・SHA256SUMS・分析md・RUN\_METADATA
を完備。床の再構成・検算は正本
\texttt{run\_N3\_N40\_stage123\_v1.py}（SHA256
\texttt{1abf2353fee2e4f56f05e7a6f149fd086885136beb61ab571b48a56b09691567}）と同一定義・物理無変更。

\end{document}
